\section{A concrete implementation plan}\label{sec:implementation}

\subsection{Priority and expected benefit}

Implement the Gram-matrix replacement first. It has the cleanest local boundary:
the public majorant theorem and its constants stay unchanged, the downstream caller
stays unchanged, and the existing Hua bound supplies all the integration machinery.
The new mathematics removes a normalization choice and its consequences.

The trace replacement is a separate, useful improvement to the theorem contract.
Its strongest established benefit is the elimination of every higher power of the
off-place operator from the comparison premise. Its net source saving is less
certain because several spectral-calculus adapters must still be formalized.
It should not be used to advertise a large measured deletion before that work is done.

\begin{center}
\begin{tabularx}{\linewidth}{@{}L{.20\linewidth}YL{.26\linewidth}@{}}
\toprule
Replacement & What disappears from the argument & Integration status\\
\midrule
Gram matrix and bordered integral & Arbitrary negative-basis reconstruction, geometric triangular normalization, and its determinant and last-entry bookkeeping & Complete mathematical adapter; not yet a compiled source replacement\\
Finite trace recovery & Reconstruction separately on generalized eigenspaces of the off-place operator, and the two-index family of moments in the caller & Complete mathematical proof; interpolation and transport support compiled\\
\bottomrule
\end{tabularx}
\end{center}

\subsection{Real-place implementation in six steps}

\begin{enumerate}
\item \textbf{Add the bordered determinant adapter.}
Use a sum index for the block decomposition of \(H\), proving
\eqref{eq:complete-square} and \eqref{eq:old-hua-adapter} together with positivity.
Reindex \(\operatorname{Fin}(i+1)\) only at the application boundary.
Factor \(S^{-1}=R^*R\) with the same positive-matrix factorization already used
in \code{HolomorphicIntegrals}. No triangular factor is required.

\item \textbf{Wrap the existing integral bound.}
Apply \code{lintegral_norm_det_add_inv_pow_le} to the parameters in
\eqref{eq:old-hua-adapter}, obtaining
\[
 \int \det\bigl(H+[U_b+A^*,v]^*[U_b+A^*,v]\bigr)^{-(k+i)}\,dA
 \leq(\det H)^{-i}(h+\norm v^2)^{-k}.
\]
Retain the source's \(k\geq7\) assumption. The exact constant
\(C_{k,i}\) need not enter the formal replacement.

\item \textbf{Keep the canonical dual frame.}
Define \(D,\omega_0,H,P,V\) by \eqref{eq:raw-data}. Reuse the pairing and
negative-space facts from the first part of \code{exists_isTriangularBasis},
stopping before the construction of \(L\).

\item \textbf{Specialize the determinant reduction.}
Apply \code{norm_det_eq_norm_det_mul_inclMatrix} directly to \(D\), as in
\eqref{eq:raw-volume}. Apply the current determinant inequality for the
Hermitian part with \(H+V^*V\) in place of \(\Id+V^*V\).
Prove the affine relation \(V_\gamma=V_{\gamma_0}+[A^*,0]\).

\item \textbf{Prove just the two invariants.}
The mixed Gram determinant inequality and Hadamard give
\((\det H)^{-1}\leq\rho^{2(i+1)}\).
The shared-vector argument and the Bessel length argument give
\(h+\norm v^2\geq\theta^{-2(j-i)}\).
These replace the determinant and final-entry estimates for \(L\).

\item \textbf{Keep the public statement.}
Replace the proof of \code{RealPlace.exists_majorant_volumeRatio_inv_pow}.
The explicit majorant is continuous, hence measurable, because its Gram determinant
is strictly positive. Its integral has exactly the old bound.
The caller \code{HasCellRank.exists_majorant_lintegral_le_kernelBound}
in \code{SchurOrbit} remains unchanged.
\end{enumerate}

For the shortest formal proof, define the majorant directly in the old
\(B_0+(RA+c_0)^*(RA+c_0)\) form, multiplied by
\((2^{i+1}/\delta)^{k+i}(\det H)^{-(k+i)}\).
Then \code{measurable_det_add_conjTranspose_mul} and the old Hua theorem apply
literally. Use the bordered determinant identity only as the bridge to the
geometric lower bound. This avoids a duplicate measurability argument.

For the determinant bound in step 5, a particularly convenient source adapter
avoids constructing the inverse Gram matrix. With
\(W=y^{-1}\omega_0\) and \(C=y^{-1}FE\), one has \(C^*W=-\Id\).
The positive Gram matrix of \([C,W]\) gives
\[
 1=\abs{\det(C^*W)}^2\leq\det(C^*C)\det H
       \leq\rho^{2(i+1)}\det H.
\]
The source already supplies \code{norm_det_sq_le_of_posSemidef} and
\code{hadamard_det_inequality}. This is the same estimate proved by projection
in Section~\ref{sec:hua}, expressed in the available determinant API.

\subsection{The exact affected regions}

The following inventory refers to the pinned revision. Counts include comments
and blank lines. Affected lines include statements that will be retained with new
proofs and mathematical roles that need shorter replacements.

\begin{center}
\begin{tabularx}{\linewidth}{@{}L{.22\linewidth}rrY@{}}
\toprule
File & Region & Lines & Role to replace\\
\midrule
\code{OneRealPlace} & 241--282 & 42 & Two reconstruction lemmas for arbitrary negative frames\\
 & 360--426 & 67 & \code{IsNegativeBasis} and its quotient theorem\\
 & 449--707 & 259 & Triangular normalization, its determinant estimates, and translation\\
 & 795--1008 & 214 & Final-column estimates and proof of the public majorant\\
\code{HermitianMatrices} & 71--99 & 29 & The triangular normalization existence lemma\\
\midrule
Total & & \textbf{611} & Gross affected source; net saving not measured\\
\bottomrule
\end{tabularx}
\end{center}

There are 22 affected declarations inside \code{OneRealPlace} and one support
declaration. The 582 affected lines inside that file are 56.5\% of its 1,030 lines.
The total of 611 lines is 0.1281\% of the full 476,849-line source tree.
This is a substantial local simplification, not evidence of a large percentage
reduction of the entire development.

The long support-lemma name
\begin{quote}\small\code{New.exists_lowerTriangular_conjTranspose_mul_mul_eq_one}\end{quote}
has only one project call site, in \code{exists_isTriangularBasis}.
The references to \code{IsNegativeBasis} and \code{IsTriangularBasis} are local to
\code{OneRealPlace}. Those facts justify the proposed removal boundary.

No complete module is shown to become unreachable. The elementary determinant
reduction, shared-vector lemma, ordinary Hua inequality, subframe determinant
formula, and unipotent-translation formula remain useful. The other Gram--Schmidt
construction in \code{Bessel/Basic}, the simultaneous-metric argument in
\code{Frames/RealPlaceGauge}, and the arithmetic module
\code{BorderedDeterminant} serve different purposes and are retained.

\subsection{Trace implementation and the formalization bottleneck}

The first port can use the simpler premise
\[
 \forall j\geq1,\quad \tr(AT^j)=\tr(A'T'^{\,j}),
\]
deferring the finite bound to a corollary. This already removes the extra index
and all powers of \(A\) above one. The archived Lean support also proves the
finite interpolation bound, so adding the finite statement afterward is straightforward
once the spectral identity is available.

\begin{enumerate}
\item Keep \code{LineRelation}, the idempotent determinant calculation, and
the local adelic raising and lowering operators. Add the spectral trace identity
of Lemma~\ref{lem:spectral-certificate} as the central finite-dimensional lemma.
\item Prove the new trace-recovery endpoint by common interpolation and positive
moment transport. If useful during migration, keep the old two-index theorem
as a wrapper: its premise supplies the new one by taking the off-place exponent
equal to one.
\item In \code{LevelCut.trace_eq_of_trace_cutLevel_eq}, specialize the action
identity to the first power of the off-place operator. The local Hecke support
and commutation proofs remain available.
\item In \code{CuspidalComparison.trace_eq_of_isCuspidalAt}, feed the original
cuspidal test function to the geometric trace comparison. The positive power
of the auxiliary operator still provides the required simple-zero condition.
The outer comparison statement need not change.
\end{enumerate}

The pinned Mathlib has the ingredients for the spectral identity, but the checked
support files do not yet assemble them. A complete formal proof needs an adapter
for \(\tr p(T)=\sum p(\lambda)\) with algebraic multiplicity, the determinant
root-product logarithm, spectral inclusion for restrictions, and the self-adjoint
decomposition of \(A\). Useful existing endpoints include
\begin{quote}\small
\code{LinearMap.trace_comp_eq_mul_of_commute_of_isNilpotent}\\
\code{LinearMap.trace_eq_sum_trace_restrict}\\
\code{Module.End.det_eq_prod_roots_charpoly_of_splits}\\
\code{LinearMap.IsSymmetric.direct_sum_isInternal}.
\end{quote}
The mathematical argument is complete, but the cost of those adapters is why a
net Lean line saving for this replacement is not asserted.

One old helper has an independent consumer:
\code{charpoly_eq_of_trace_pow_eq} is used by
\code{EqualMultiplicities.charpoly_heckeOnGood_eq}.
It must remain, or receive a separate replacement. The new trace theorem does not
make all generalized-eigenspace or characteristic-polynomial infrastructure dead code.

\subsection{Acceptance criteria for an integrated change}

For the real-place replacement, the decisive test is that the original public
majorant statement and the unchanged \code{SchurOrbit} caller compile, including
the empty-column case \(i=0\). For trace recovery, the new theorem must retain both
commutation hypotheses and allow nonnormal raising operators; the repeated-root
example in Section~\ref{sec:trace} is a useful guard against an accidental
diagonalizability assumption.

After those local ports, rebuild \code{Main} at the pinned dependencies and inspect
the public theorem assumptions. Then measure the source delta including new helpers.
That would establish the implementation and its net saving. Those integration checks
have not been performed for this proposal archive.
